\section{Turbulence Closure Models}
\label{sec:closures}

We evaluate nine turbulence closure models spanning the full hierarchy from zero-equation algebraic to random forest.
All models are implemented in the same solver and share the same spatial discretization, time integration, and Poisson solver.
Table~\ref{tab:model_summary} summarizes their key properties.

\begin{table}[htbp]
\centering
\caption{Summary of turbulence closure models evaluated. FLOPs/cell counts are for the turbulence update only (excluding shared solver operations). Weight sizes refer to the model parameters stored on device.}
\label{tab:model_summary}
\begin{tabular}{lllrrl}
\toprule
Model & Type & Architecture & Parameters & FLOPs/cell & Weights \\
\midrule
Baseline       & Algebraic    & Mixing length              & 0       & $\sim$2    & --- \\
$k$-$\omega$   & Transport    & 2-equation                 & 0       & $\sim$50   & --- \\
SST            & Transport    & 2-equation + blending      & 0       & $\sim$80   & --- \\
EARSM          & Algebraic RS & SST + tensor algebra       & 0       & $\sim$120  & --- \\
GEP            & Algebraic RS & Symbolic regression        & 0       & $\sim$100  & --- \\
MLP            & NN (scalar)  & $5 \to 32 \to 32 \to 1$    & 1{,}249  & $\sim$1{,}249 & \SI{56}{KB} \\
MLP-Large      & NN (scalar)  & $5 \to 128^4 \to 1$        & 50{,}049 & $\sim$50{,}049 & \SI{896}{KB} \\
TBNN           & NN (tensor)  & $5 \to 64^3 \to 10$        & 9{,}354  & $\sim$10{,}300 & \SI{196}{KB} \\
PI-TBNN        & NN (tensor)  & Same as TBNN               & 9{,}354  & $\sim$10{,}300 & \SI{196}{KB} \\
TBRF           & RF (tensor)  & 10 coeff.\ $\times$ 10 trees & $\sim$2.8M nodes & $\sim$2{,}000 branches & \SI{56}{MB} \\
\bottomrule
\end{tabular}
\end{table}

\subsection{Classical models}

\subsubsection{Baseline mixing length}

The algebraic baseline computes the eddy viscosity from a mixing length $\ell_{\mathrm{mix}}$ and the local velocity gradient:
\begin{equation}
    \nut = \ell_{\mathrm{mix}}^2 \left| \frac{\partial u}{\partial y} \right|.
\end{equation}
This requires one division and one multiply per cell, reading 3 velocity fields and writing 1 eddy viscosity field.
On GPU, this is a memory-bandwidth-bound operation with perfect parallelism.

\subsubsection{$k$-$\omega$ transport model}

The $k$-$\omega$ model \citep{Wilcox1988} solves transport equations for the turbulent kinetic energy $k$ and the specific dissipation rate $\omega$:
\begin{align}
    \frac{\partial k}{\partial t} + u_j \frac{\partial k}{\partial x_j} &= P_k - \beta^* k \omega + \frac{\partial}{\partial x_j}\left[(\nu + \sigma_k \nut)\frac{\partial k}{\partial x_j}\right], \\
    \frac{\partial \omega}{\partial t} + u_j \frac{\partial \omega}{\partial x_j} &= \alpha \frac{\omega}{k} P_k - \beta \omega^2 + \frac{\partial}{\partial x_j}\left[(\nu + \sigma_\omega \nut)\frac{\partial \omega}{\partial x_j}\right].
\end{align}
Each equation involves convective and diffusive flux computation (stencil operations) plus source and sink terms, totaling approximately 50~FLOPs per cell.
Destruction terms ($\beta^* k \omega$, $\beta \omega^2$) are treated point-implicitly to avoid stiffness-induced blow-up at wall cells, where $\omega_{\mathrm{wall}} \sim O(10^5)$.

\subsubsection{SST $k$-$\omega$}

The SST model \citep{Menter1994} augments $k$-$\omega$ with blending functions $F_1$ and $F_2$ that transition between near-wall ($k$-$\omega$) and free-stream ($k$-$\varepsilon$) formulations, plus a cross-diffusion term $\mathrm{CD}_{k\omega}$.
Each blending function requires wall distance, $\nabla k \cdot \nabla \omega$, and several auxiliary quantities, raising the cost to approximately 80~FLOPs per cell (2.5$\times$ the base $k$-$\omega$ model).

\subsubsection{EARSM}

The explicit algebraic Reynolds stress model extends SST by computing the anisotropy tensor $a_{ij}$ from an algebraic relation involving the normalized strain rate $\Shat_{ij}$ and rotation rate $\Omhat_{ij}$ tensors:
\begin{equation}
    a_{ij} = f(\Shat, \Omhat),
\end{equation}
where the functional form follows Pope, Wallin--Johansson, or Gatski--Speziale closures.
The tensor algebra (3$\times$3 matrix multiplications, traces, invariant computation) adds approximately 40~FLOPs per cell on top of the SST transport, totaling $\sim$120~FLOPs per cell.
Despite the higher per-cell cost, EARSM retains the stencil-based memory access pattern of transport models and remains memory-bandwidth-bound on GPU.

\subsubsection{GEP}

The gene expression programming model \citep{Weatheritt2016} uses a symbolic expression discovered via genetic programming to compute the Reynolds stress anisotropy.
The expression is algebraic (no transport equations beyond SST), with a cost of $\sim$100~FLOPs per cell.

\subsection{Neural network architectures}

All neural network models take the same 5 scalar invariants of the non-dimensionalized strain rate and rotation rate tensors as input (Section~\ref{sec:training}).
The architectures differ in capacity, output dimensionality, and how the output maps to a turbulence closure.

\subsubsection{MLP (scalar closure)}

The multilayer perceptron predicts a scalar proxy for the eddy viscosity magnitude:
\begin{equation}
    |b| = f_{\mathrm{MLP}}(\hat{\lambda}_1, \ldots, \hat{\lambda}_5),
\end{equation}
with architecture $5 \to 32 \to 32 \to 1$ (1{,}249 parameters).
Hidden layers use $\tanh$ activations; the output layer is linear.
The inference cost per cell is:
\begin{itemize}
    \item Layer 0 ($5 \to 32$): 160 multiplies + 32 adds + 32 $\tanh$
    \item Layer 1 ($32 \to 32$): 1{,}024 multiplies + 32 adds + 32 $\tanh$
    \item Layer 2 ($32 \to 1$): 32 multiplies + 1 add
    \item \textbf{Total}: 1{,}249 FLOPs + 64 $\tanh$ evaluations
\end{itemize}
The weight matrix (\SI{56}{KB}) fits in GPU L1 cache, enabling broadcast to all threads with minimal memory traffic.

\subsubsection{MLP-Large (scalar closure)}

A larger MLP with architecture $5 \to 128 \to 128 \to 128 \to 128 \to 1$ (50{,}049 parameters) serves as a capacity scaling test.
The dominant cost is four $128 \times 128$ matrix-vector products (16{,}384 multiplies each), totaling $\sim$50{,}000~FLOPs per cell with 512~$\tanh$ evaluations.
The weight matrix (\SI{896}{KB}) exceeds L1 but fits in L2 cache.

\subsubsection{TBNN (tensor basis closure)}
\label{sec:tbnn}

The tensor basis neural network \citep{Ling2016} predicts 10 scalar coefficients $g^{(n)}$ that are contracted with the Pope integrity basis \citep{Pope1975} to reconstruct the anisotropy tensor:
\begin{equation}
    b_{ij} = \sum_{n=1}^{10} g^{(n)}(\hat{\lambda}_1, \ldots, \hat{\lambda}_5) \, T^{(n)}_{ij},
    \label{eq:tbnn}
\end{equation}
where $T^{(n)}$ are the 10 symmetric, traceless tensor basis functions of $\Shat$ and $\Omhat$.
The architecture is $5 \to 64 \to 64 \to 64 \to 10$ (9{,}354 parameters) with $\tanh$ hidden activations.

The per-cell computational cost breaks down as follows:
\begin{itemize}
    \item \textbf{Feature computation}: velocity gradients (stencil), $S_{ij}$ and $\Omega_{ij}$ (symmetrize/antisymmetrize), non-dimensionalize by $k/\varepsilon$, compute 5 invariants (matrix products + traces)---shared with EARSM, $\sim$80 FLOPs.
    \item \textbf{Tensor basis computation}: 10 basis tensors $T^{(n)}$ from products of $\Shat$ and $\Omhat$, each requiring one or more $3 \times 3$ matrix multiplies---$\sim$900~FLOPs.
    \item \textbf{NN forward pass}: 4 layers totaling 9{,}354 FLOPs + 192 $\tanh$ evaluations. The cost is dominated by the two $64 \times 64$ hidden layers (4{,}096 FLOPs each, comprising 87\% of the network's compute).
    \item \textbf{Tensor contraction}: $b_{ij} = \sum_n g^{(n)} T^{(n)}_{ij}$---60 multiplies + 60 adds.
    \item \textbf{Total}: $\sim$10{,}300 FLOPs/cell, approximately 7.5$\times$ the MLP.
\end{itemize}

The TBNN weight matrix (\SI{196}{KB}) fits in GPU L2 cache.
The architecture guarantees Galilean invariance, symmetry, and zero trace of the predicted $b_{ij}$ by construction.

\subsubsection{PI-TBNN (physics-informed TBNN)}

The PI-TBNN shares the same architecture as the TBNN but is trained with an augmented loss function that penalizes violations of the realizability constraints $b_{ii} \geq -1/3$ and $b_{ii} \leq 2/3$ (see Section~\ref{sec:training}).
At inference time, the PI-TBNN is computationally identical to the TBNN.

\subsubsection{TBRF (tensor basis random forest)}

The tensor basis random forest \citep{Kaandorp2020} replaces the neural network with an ensemble of decision trees.
For each of the 10 basis coefficients $g^{(n)}$, a separate random forest of 10--200 trees (depth 20) is trained on least-squares-extracted coefficients.

The inference cost per cell differs fundamentally from neural networks:
\begin{itemize}
    \item For each cell, for each of 10 coefficients, traverse 10 trees to depth $\sim$20: this requires $10 \times 10 \times 20 = 2{,}000$ comparison-and-branch operations.
    \item Each branch decision accesses a different memory location in the tree structure (feature index, threshold, child pointers), producing scattered, non-coalesced memory accesses.
    \item The 10-tree TBRF requires \SI{56}{MB} of tree data, which does not fit in any GPU cache level (L1: \SI{128}{KB}, L2: \SI{48}{MB} on A100).
    \item Warp divergence is severe: cells in the same warp follow different paths through each tree, serializing execution.
\end{itemize}

Contrast this with neural network inference, where every cell performs the \emph{same} matrix-vector multiply with the \emph{same} weights---a perfectly regular, data-parallel operation with high arithmetic intensity and coalesced memory access.

\subsection{Computational anatomy: why NNs are expensive on GPU}
\label{sec:comp_anatomy}

The fundamental mismatch between neural network inference and GPU architecture manifests in three ways:

\paragraph{Layer-sequential execution.}
Each MLP layer depends on the previous layer's output.
For a TBNN with 4 layers, this means 4 serial phases within each kernel.
Parallelism is only over (cell, neuron) pairs; the depth dimension is sequential.
This limits the arithmetic intensity compared to stencil operations that process all cells independently.

\paragraph{Activation functions.}
The $\tanh$ function costs approximately 20 GPU cycles per evaluation versus $\sim$4 cycles for a multiply.
With 192 $\tanh$ evaluations per cell for the TBNN, activation functions account for a significant fraction of the total compute.

\paragraph{Memory access patterns.}
Classical stencil-based closures (transport models) access compact neighborhoods of large arrays with perfect spatial locality and coalesced GPU memory access.
Neural network weights must be broadcast to all threads---this is efficient for small models (MLP, \SI{56}{KB}, fits in L1) but increasingly costly as model size grows.
The TBRF's tree data (\SI{56}{MB}) causes cache thrashing and random memory accesses, making it fundamentally poorly suited for GPU execution.
